%% kapa: exam=1 round=27 year=2016 period=1 form=A subject=civil_law subject_ko=민법 seq=23 number=23
%% source: Q-Net 감정평가사 기출문제 (cst00309) artlSeq=5169402 file=extracted/1008/5169402/제27회 감정평가사 1차 문제/제27회 감정평가사 1교시문제(A형).pdf
%% conversion: kordoc 4.24.51 -> tools/gen1.py (auto); review=proofread-sonnet; flags=none
\begin{question}{23}
선의취득에 관한 설명으로 옳지 않은 것은? (다툼이 있으면 판례에 따름)
\begin{choices}
  \item 경매에 의한 동산의 취득에는 선의취득이 인정되지 않는다.
  \item 점유개정은 선의취득에서의 인도의 방법으로 인정되지 않는다.
  \item 간이인도는 선의취득에서의 인도의 방법으로 인정된다.
  \item 저당권은 선의취득의 대상이 아니지만 동산질권은 선의취득의 대상이 된다.
  \item 물권적 합의가 동산의 인도보다 먼저 행하여지면 양수인의 선의ㆍ무과실은 인도된 때를 기준으로 판단한다.
\end{choices}
\end{question}
\answer{1}
